% ============================================================
% Question Bank: ch02-09 Financial Literacy: Budgeting, Saving, and Investing
% Topic Title: Financial Literacy: Budgeting, Saving, and Investing
% CCSS: Supplementary
% Questions: 27 (15 MC + 6 SA + 6 Graphical)
% ============================================================

% --- Multiple Choice Questions (15) ---

\begin{practiceQuestion}{2-9-q01}{mc}
\begin{questionText}
A family earns $\$4{,}000$ per month and spends $\$3{,}400$. What percent of their income do they save?
\end{questionText}
\choiceA{$10\%$}
\choiceB{$15\%$}
\choiceC{$17\%$}
\choiceD{$85\%$}
\correctAnswer{B}
\explanation{Savings: $4{,}000 - 3{,}400 = \$600$. Rate: $600 \div 4{,}000 = 0.15 = 15\%$.}
\end{practiceQuestion}

\begin{practiceQuestion}{2-9-q02}{mc}
\begin{questionText}
Which of the following is a \textbf{fixed} expense?
\end{questionText}
\choiceA{Groceries}
\choiceB{Monthly rent}
\choiceC{Entertainment}
\choiceD{Clothing}
\correctAnswer{B}
\explanation{Rent stays the same amount each month, making it a fixed expense. Groceries, entertainment, and clothing vary month to month.}
\end{practiceQuestion}

\begin{practiceQuestion}{2-9-q03}{mc}
\begin{questionText}
A student's monthly budget is $\$300$. She spends $\$120$ on food, $\$60$ on transportation, and $\$45$ on entertainment. How much is left for savings?
\end{questionText}
\choiceA{$\$25$}
\choiceB{$\$45$}
\choiceC{$\$75$}
\choiceD{$\$225$}
\correctAnswer{C}
\explanation{Total spent: $120 + 60 + 45 = \$225$. Left: $300 - 225 = \$75$.}
\end{practiceQuestion}

\begin{practiceQuestion}{2-9-q04}{mc}
\begin{questionText}
Which of the following best describes the difference between saving and investing?
\end{questionText}
\choiceA{Saving has higher risk and higher return}
\choiceB{Investing has lower risk and lower return}
\choiceC{Saving is lower risk; investing may offer higher returns but carries more risk}
\choiceD{There is no difference}
\correctAnswer{C}
\explanation{Saving (bank accounts) is low risk with modest returns. Investing (stocks, bonds) carries more risk but has the potential for higher returns over time.}
\end{practiceQuestion}

\begin{practiceQuestion}{2-9-q05}{mc}
\begin{questionText}
Alex earns $\$2{,}500$ per month. Financial advisors recommend saving at least $20\%$ of income. How much should Alex save each month?
\end{questionText}
\choiceA{$\$200$}
\choiceB{$\$250$}
\choiceC{$\$375$}
\choiceD{$\$500$}
\correctAnswer{D}
\explanation{$20\%$ of $\$2{,}500 = 0.20 \times 2{,}500 = \$500$.}
\end{practiceQuestion}

\begin{practiceQuestion}{2-9-q06}{mc}
\begin{questionText}
A family's monthly budget allocates $30\%$ to housing, $15\%$ to food, $10\%$ to transportation, $10\%$ to savings, and the rest to other expenses. If income is $\$5{,}000$, how much goes to ``other expenses''?
\end{questionText}
\choiceA{$\$500$}
\choiceB{$\$1{,}000$}
\choiceC{$\$1{,}500$}
\choiceD{$\$1{,}750$}
\correctAnswer{D}
\explanation{Allocated: $30 + 15 + 10 + 10 = 65\%$. Other: $100\% - 65\% = 35\%$. Amount: $0.35 \times 5{,}000 = \$1{,}750$.}
\end{practiceQuestion}

\begin{practiceQuestion}{2-9-q07}{mc}
\begin{questionText}
A teenager saves $\$25$ per week. How many weeks will it take to save $\$400$ for a new phone?
\end{questionText}
\choiceA{$10$ weeks}
\choiceB{$14$ weeks}
\choiceC{$16$ weeks}
\choiceD{$20$ weeks}
\correctAnswer{C}
\explanation{$400 \div 25 = 16$ weeks.}
\end{practiceQuestion}

\begin{practiceQuestion}{2-9-q08}{mc}
\begin{questionText}
Which of the following is a \textbf{variable} expense?
\end{questionText}
\choiceA{Car payment}
\choiceB{Mortgage}
\choiceC{Electric bill}
\choiceD{Insurance premium}
\correctAnswer{C}
\explanation{Electric bills change month to month based on usage, making them a variable expense. Car payments, mortgages, and insurance premiums are typically fixed.}
\end{practiceQuestion}

\begin{practiceQuestion}{2-9-q09}{mc}
\begin{questionText}
A family reduces their variable spending by $12\%$. If they were spending $\$750$ per month on variables, how much do they save per month?
\end{questionText}
\choiceA{$\$12$}
\choiceB{$\$75$}
\choiceC{$\$90$}
\choiceD{$\$120$}
\correctAnswer{C}
\explanation{$750 \times 0.12 = \$90$ per month saved.}
\end{practiceQuestion}

\begin{practiceQuestion}{2-9-q10}{mc}
\begin{questionText}
Tina invests $\$500$ in a savings account that earns $3\%$ per year. Marcus invests $\$500$ in an index fund that earns $8\%$ per year. After $1$ year (simple interest), how much more does Marcus earn?
\end{questionText}
\choiceA{$\$5$}
\choiceB{$\$15$}
\choiceC{$\$25$}
\choiceD{$\$40$}
\correctAnswer{C}
\explanation{Tina: $500 \times 0.03 = \$15$. Marcus: $500 \times 0.08 = \$40$. Difference: $40 - 15 = \$25$.}
\end{practiceQuestion}

\begin{practiceQuestion}{2-9-q11}{mc}
\begin{questionText}
What does the equation \textbf{Income $-$ Expenses $=$ Savings} tell you?
\end{questionText}
\choiceA{How much you owe}
\choiceB{How much money is left over after spending}
\choiceC{How much you should invest}
\choiceD{How much tax you pay}
\correctAnswer{B}
\explanation{The equation shows that whatever you earn minus what you spend is what you save. Savings is what remains after expenses.}
\end{practiceQuestion}

\begin{practiceQuestion}{2-9-q12}{mc}
\begin{questionText}
A budget shows $\$1{,}200$ for fixed expenses and $\$600$ for variable expenses. The total income is $\$2{,}400$. What savings rate does this budget achieve?
\end{questionText}
\choiceA{$10\%$}
\choiceB{$20\%$}
\choiceC{$25\%$}
\choiceD{$50\%$}
\correctAnswer{C}
\explanation{Total expenses: $1{,}200 + 600 = \$1{,}800$. Savings: $2{,}400 - 1{,}800 = \$600$. Rate: $600 \div 2{,}400 = 0.25 = 25\%$.}
\end{practiceQuestion}

\begin{practiceQuestion}{2-9-q13}{mc}
\begin{questionText}
If you save $\$150$ per month, how much will you save in $2$ years?
\end{questionText}
\choiceA{$\$1{,}800$}
\choiceB{$\$3{,}000$}
\choiceC{$\$3{,}600$}
\choiceD{$\$4{,}200$}
\correctAnswer{C}
\explanation{$2$ years $= 24$ months. $150 \times 24 = \$3{,}600$.}
\end{practiceQuestion}

\begin{practiceQuestion}{2-9-q14}{mc}
\begin{questionText}
Which is the best reason to have an emergency fund?
\end{questionText}
\choiceA{To invest in stocks}
\choiceB{To pay for daily groceries}
\choiceC{To cover unexpected expenses like medical bills or car repairs}
\choiceD{To buy holiday gifts}
\correctAnswer{C}
\explanation{An emergency fund covers unexpected costs so you do not need to borrow money or use credit cards for surprises.}
\end{practiceQuestion}

\begin{practiceQuestion}{2-9-q15}{mc}
\begin{questionText}
A worker earns $\$3{,}200$ per month. Her fixed expenses are $\$1{,}800$ and variable expenses are $\$900$. She wants to save at least $10\%$ of income. Does she meet this goal?
\end{questionText}
\choiceA{Yes --- she saves $\$500$ ($15.6\%$)}
\choiceB{No --- she only saves $\$200$ ($6.25\%$)}
\choiceC{Yes --- she saves $\$320$ ($10\%$)}
\choiceD{No --- she only saves $\$300$ ($9.4\%$)}
\correctAnswer{A}
\explanation{Savings: $3{,}200 - 1{,}800 - 900 = \$500$. Rate: $500 \div 3{,}200 = 0.15625 \approx 15.6\%$. This exceeds $10\%$.}
\end{practiceQuestion}

% --- Short Answer Questions (6) ---

\begin{practiceQuestion}{2-9-q16}{sa}
\begin{questionText}
A family earns $\$5{,}500$ per month and saves $\$825$. What percent of their income do they save?
\end{questionText}
\correctAnswer{$15\%$}
\explanation{$825 \div 5{,}500 = 0.15 = 15\%$.}
\end{practiceQuestion}

\begin{practiceQuestion}{2-9-q17}{sa}
\begin{questionText}
A student wants to save $\$1{,}200$ in $8$ months. How much must she save per month?
\end{questionText}
\correctAnswer{$\$150$}
\explanation{$1{,}200 \div 8 = \$150$ per month.}
\end{practiceQuestion}

\begin{practiceQuestion}{2-9-q18}{sa}
\begin{questionText}
Monthly income is $\$3{,}000$. Fixed expenses are $\$1{,}500$ and variable expenses are $\$800$. What is the monthly savings?
\end{questionText}
\correctAnswer{$\$700$}
\explanation{$3{,}000 - 1{,}500 - 800 = \$700$.}
\end{practiceQuestion}

\begin{practiceQuestion}{2-9-q19}{sa}
\begin{questionText}
A family decides to cut $15\%$ from their $\$600$ monthly food budget. How much will they spend on food?
\end{questionText}
\correctAnswer{$\$510$}
\explanation{Cut: $600 \times 0.15 = \$90$. New amount: $600 - 90 = \$510$.}
\end{practiceQuestion}

\begin{practiceQuestion}{2-9-q20}{sa}
\begin{questionText}
If you earn $\$2{,}000$ per month and your savings rate is $12\%$, how much do you save in $6$ months?
\end{questionText}
\correctAnswer{$\$1{,}440$}
\explanation{Monthly savings: $2{,}000 \times 0.12 = \$240$. In $6$ months: $240 \times 6 = \$1{,}440$.}
\end{practiceQuestion}

\begin{practiceQuestion}{2-9-q21}{sa}
\begin{questionText}
An investment earns $\$160$ in one year on a $\$2{,}000$ investment. What is the annual return rate?
\end{questionText}
\correctAnswer{$8\%$}
\explanation{$160 \div 2{,}000 = 0.08 = 8\%$.}
\end{practiceQuestion}

% --- Graphical Questions (3 GMC + 3 GSA) ---

\begin{practiceQuestion}{2-9-q22}{gmc}
\begin{questionText}
The pie chart shows a family's monthly budget. If their income is $\$4{,}000$, how much goes to housing?

\begin{center}
\begin{tikzpicture}[scale=2.1]
  \fill[funBlue!50] (0,0) -- (90:1) arc (90:198:1) -- cycle;
  \fill[funGreen!50] (0,0) -- (198:1) arc (198:270:1) -- cycle;
  \fill[funOrange!50] (0,0) -- (270:1) arc (270:342:1) -- cycle;
  \fill[funRed!50] (0,0) -- (342:1) arc (342:450:1) -- cycle;
  \draw (0,0) circle (1);
  \draw (0,0) -- (90:1); \draw (0,0) -- (198:1); \draw (0,0) -- (270:1); \draw (0,0) -- (342:1);
  \node[font=\footnotesize,align=center] at (144:0.6) {Housing\\$30\%$};
  \node[font=\footnotesize,align=center] at (234:0.55) {Food\\$20\%$};
  \node[font=\footnotesize,align=center] at (306:0.58) {Transport\\$20\%$};
  \node[font=\footnotesize,align=center] at (36:0.6) {Other\\$30\%$};
\end{tikzpicture}
\end{center}
\end{questionText}
\choiceA{$\$800$}
\choiceB{$\$1{,}000$}
\choiceC{$\$1{,}200$}
\choiceD{$\$1{,}500$}
\correctAnswer{C}
\explanation{Housing: $30\%$ of $\$4{,}000 = 0.30 \times 4{,}000 = \$1{,}200$.}
\end{practiceQuestion}

\begin{practiceQuestion}{2-9-q23}{gmc}
\begin{questionText}
The bar graph shows a student's spending by category over one month. If her total income was $\$500$, what percent went to savings?

\begin{center}
\begin{tikzpicture}[scale=0.5]
  \draw[thick,->] (0,0) -- (0,6) node[above]{\small Amount (\$)};
  \draw[thick,->] (0,0) -- (10,0);
  \fill[funBlue!60] (0.5,0) rectangle (2,4);
  \fill[funGreen!60] (2.5,0) rectangle (4,3);
  \fill[funOrange!60] (4.5,0) rectangle (6,1.5);
  \fill[funRed!60] (6.5,0) rectangle (8,1.5);
  \node[below,font=\small] at (1.25,0) {Food};
  \node[below,font=\small] at (3.25,0) {Fun};
  \node[below,font=\small] at (5.25,0) {Trans.};
  \node[below,font=\small] at (7.25,0) {Other};
  \foreach \y/\lab in {1/50,2/100,3/150,4/200,5/250} {
    \draw (-0.1,\y) -- (0.1,\y) node[left,font=\small,xshift=-2pt] {\lab};
  }
  \node[above,font=\small] at (1.25,4) {$\$200$};
  \node[above,font=\small] at (3.25,3) {$\$150$};
  \node[above,font=\small] at (5.25,1.5) {$\$75$};
  \node[above,font=\small] at (7.25,1.5) {$\$25$};
\end{tikzpicture}
\end{center}
\end{questionText}
\choiceA{$5\%$}
\choiceB{$10\%$}
\choiceC{$15\%$}
\choiceD{$20\%$}
\correctAnswer{B}
\explanation{Total spent: $200 + 150 + 75 + 25 = \$450$. Savings: $500 - 450 = \$50$. Rate: $50 \div 500 = 0.10 = 10\%$.}
\end{practiceQuestion}

\begin{practiceQuestion}{2-9-q24}{gmc}
\begin{questionText}
The table shows two investment options for $\$1{,}000$ over $1$ year (simple interest). Which earns more and by how much?

\begin{center}
\begin{tabular}{|l|c|c|}
\hline
\textbf{Option} & \textbf{Type} & \textbf{Annual Return} \\
\hline
Savings Account & Low risk & $2\%$ \\
\hline
Stock Fund & Higher risk & $9\%$ \\
\hline
\end{tabular}
\end{center}
\end{questionText}
\choiceA{Savings by $\$20$}
\choiceB{Stock fund by $\$70$}
\choiceC{Stock fund by $\$90$}
\choiceD{They earn the same}
\correctAnswer{B}
\explanation{Savings: $1{,}000 \times 0.02 = \$20$. Stock: $1{,}000 \times 0.09 = \$90$. Difference: $90 - 20 = \$70$.}
\end{practiceQuestion}

\begin{practiceQuestion}{2-9-q25}{gsa}
\begin{questionText}
The table shows a household's monthly expense breakdown. If income is $\$3{,}500$, how much is saved?

\begin{center}
\begin{tabular}{|l|c|}
\hline
\textbf{Category} & \textbf{Percent} \\
\hline
Housing & $35\%$ \\
\hline
Food & $20\%$ \\
\hline
Transportation & $10\%$ \\
\hline
Utilities & $8\%$ \\
\hline
Entertainment & $7\%$ \\
\hline
Savings & ? \\
\hline
\end{tabular}
\end{center}
\end{questionText}
\correctAnswer{$\$700$}
\explanation{Allocated: $35 + 20 + 10 + 8 + 7 = 80\%$. Savings: $100\% - 80\% = 20\%$. Amount: $0.20 \times 3{,}500 = \$700$.}
\end{practiceQuestion}

\begin{practiceQuestion}{2-9-q26}{gsa}
\begin{questionText}
The line graph shows a student's savings balance over $5$ months. How much did the student save per month on average?

\begin{center}
\begin{tikzpicture}[scale=0.7]
  \draw[thick,->] (0,0) -- (6,0) node[right]{\small Month};
  \draw[thick,->] (0,0) -- (0,5.5) node[above]{\small Balance (\$)};
  \foreach \x/\lab in {0/0,1/1,2/2,3/3,4/4,5/5} {
    \draw (\x,-0.1) -- (\x,0.1) node[below=3pt,font=\small] {\lab};
  }
  \foreach \y/\lab in {1/100,2/200,3/300,4/400,5/500} {
    \draw (-0.1,\y) -- (0.1,\y) node[left,font=\small,xshift=-2pt] {\lab};
  }
  \draw[thick,funGreen,mark=*,mark size=2pt] plot coordinates {(0,0.5) (1,1.2) (2,2) (3,2.8) (4,3.5) (5,4.5)};
  \node[right,font=\small,funGreen] at (5,4.5) {$\$450$};
  \node[left,font=\small,funGreen] at (0,0.5) {$\$50$};
\end{tikzpicture}
\end{center}
\end{questionText}
\correctAnswer{$\$80$}
\explanation{Net savings: $450 - 50 = \$400$ over $5$ months. Average: $400 \div 5 = \$80$ per month.}
\end{practiceQuestion}

\begin{practiceQuestion}{2-9-q27}{gsa}
\begin{questionText}
The table compares two families' budgets. Which family has the higher savings rate, and what is it?

\begin{center}
\begin{tabular}{|l|c|c|c|}
\hline
 & \textbf{Income} & \textbf{Fixed} & \textbf{Variable} \\
\hline
Family A & $\$4{,}000$ & $\$2{,}200$ & $\$1{,}200$ \\
\hline
Family B & $\$6{,}000$ & $\$3{,}500$ & $\$1{,}800$ \\
\hline
\end{tabular}
\end{center}
\end{questionText}
\correctAnswer{Family A ($15\%$)}
\explanation{A: Savings $= 4{,}000 - 2{,}200 - 1{,}200 = \$600$. Rate $= 600 \div 4{,}000 = 15\%$. B: Savings $= 6{,}000 - 3{,}500 - 1{,}800 = \$700$. Rate $= 700 \div 6{,}000 \approx 11.7\%$. Family A has the higher rate.}
\end{practiceQuestion}
